\section{Sharp comparison at depth two}

A star consists of a root with its pairwise disjoint leaves. For one
such component write $r$ for the root weight, $b=\norm{q}$ for the leaf
norm and $A=\E|\sum_i q_i\epsilon_i|$ for the first absolute moment of
the independent fair leaf signs. The recursions give
\[
 \Vstar=r^2+b^2+2rA,\qquad U=(r+b)^2.
\]

\begin{theorem}[Depth two]\label{thm:depth2}
For every finite laminar forest of depth at most two,
$\Vstar\ge c_2U$ with sharp constant $c_2=(1+1/\sqrt2)/2$.
If $U>0$, equality holds precisely when each positive-weight
component has exactly two equal nonzero leaf magnitudes and root
magnitude equal to its leaf norm.
\end{theorem}
\begin{proof}
Subtract $c_2U$ on each component:
\[
 \Vstar-c_2U=(1-c_2)(r-b)^2+2r(A-b/\sqrt2).
\]
The real sharp $L^1$ Khintchine inequality gives $A\ge b/\sqrt2$
\cite{szarek,haagerup}. For $b>0$, equality in that inequality requires
exactly two equal nonzero magnitudes \cite{de}. A positive-weight
component with zero deficit therefore has $r=b>0$ and the described
leaves. The two forest objectives add across maximal components.
For sharpness take the root $(0,2]$, leaves $(0,1],(1,2]$, and
weights $(\sqrt2,1,1)$. Then $A=1$, $\Vstar=4+2\sqrt2$, and $U=8$.
At $r=b>0$ the comparison reduces to the sharp Khintchine bound.
\end{proof}

\section{An outer interval with two or three small stars}

Consider one outer root of magnitude $r$ whose children are stars
with at most two nonzero leaf weights each. For a child with root
magnitude $m$ and leaf magnitudes $p\ge q\ge0$, the canonical total in
Theorem~\ref{thm:tree} has law $(m+p)\epsilon+q\eta$ for independent fair
signs. For $p>q$ the canonical root copies the heavier sign. For $p=q$
copying that sign gives the same distribution as the fair tie rule.
When $p=q>0$, the two cancelling leaf configurations are equiprobable,
and the copied sign takes opposite values on them. If $p=q=0$,
a fresh fair sign gives total $m\epsilon$. Zero or absent leaves may
be padded algebraically. The corresponding child Gram norm is
$m+\sqrt{p^2+q^2}$.

At fixed effective pair $(a,b)=(m+p,q)$ we have $a\ge p\ge b\ge0$.
The Gram norm $a-p+\sqrt{p^2+b^2}$ is largest at $p=b$, strictly so
for $b>0$ and $p>b$. Thus the replacement $(m,p,q)=(a-b,b,b)$
preserves the canonical inner law and weakly increases $U$. Set
$k=\sqrt2-1$. In the reduced family, independently over the effective
pairs $(a_j,b_j)$, put
\begin{align*}
 S&=\sum_j(a_j\epsilon_j+b_j\eta_j), & A&=\E|S|,\\
 Q&=\sum_j(a_j^2+b_j^2), & B&=\sqrt{\sum_j(a_j+kb_j)^2}.
\end{align*}
The quotient is
\begin{equation}
 \frac{\Vstar}{U}=\frac{r^2+2rA+Q}{(r+B)^2}.
 \label{eq:ratio}
\end{equation}
A reduced bound applies to the original problem because its numerator is
unchanged and its denominator can only increase. For six unit inner
weights and $r=1$, the exact values are $A=5/2$, $Q=10$ and
$B=2+\sqrt2$. Thus the quotient is $16/(11+6\sqrt2)<c_2$.

\begin{lemma}[Dominant effective coefficient]\label{lem:dominant}
If the largest effective coefficient is at least the sum of all others
and $U>0$, the quotient in~\eqref{eq:ratio} is strictly above
$\gamma=3/(15-8\sqrt2)$.
\end{lemma}
\begin{proof}
For a nonzero effective vector let $a,b$ be its two largest coefficients,
padding zeros when needed. Conditioning on the other signs gives $A=a$.
To upper-bound $B$ we may pair $a$ with $b$. Indeed, if their partners
were $c,d$, replacing $(a,c),(b,d)$ by $(a,b),(c,d)$ increases the
squared norm by $2k(b-c)(a-b-c-d)\ge0$ before sorting the last pair.
Sorting that pair can only increase it. Let $L$ sum the coefficients
other than $a,b$, $e=a+kb$, and $J$ be the remaining pair norm. Since
$(x+ky)^2\le x(x+y)$ for $x\ge y\ge0$, $J^2\le bL$ and
$e\ge(1+k)b+L$. If $J>0$,
\[
 \frac{J}{\sqrt{e^2+J^2}+e}
 \le \frac{\sqrt{bL}}{2((1+k)b+L)}
 \le \theta:=\frac1{4\sqrt{1+k}} .
\]
Thus $B\le a+kb+\theta J$, also for $J=0$.
Cauchy--Schwarz gives $J^2\le(1+k^2)Q_{\rm rest}$ for the remaining
squared coefficients $Q_{\rm rest}$. Apply it again to
$(r+a,b,\sqrt{Q_{\rm rest}})$ to bound~\eqref{eq:ratio} below by
\[
 \frac1{(1+k^2)(1+\theta^2)}
 =\frac{8\sqrt2+8}{16\sqrt2+1}>\gamma .
\]
The strict comparison follows from $8\sqrt2>11$. If the effective vector
is zero, the nonzero root alone gives quotient one.
\end{proof}

\begin{theorem}[Seven weights]\label{thm:seven}
For one outer root and at most two disjoint child stars with at most
two active leaves each, $\Vstar/U\ge\gamma$ whenever $U>0$, where
\[
 \gamma=\frac3{15-8\sqrt2}.
\]
This is the minimum over the class of such trees and weights. Equality
requires two child stars with two leaves available in each. Up to
common positive scaling and relabeling, the equality magnitudes are
outer weight $2$, both child root weights $(3\sqrt2-1)/2$, and four
unit leaves. Every other leaf has weight zero. Signs are arbitrary.
\end{theorem}
\begin{proof}
Sort the four effective coefficients $a\ge b\ge c\ge d\ge0$.
The numerator is invariant under pairings. For two pairs denote the
squared norm by $B_{ab,cd}^2=(a+kb)^2+(c+kd)^2$ and use analogous
notation for other pairings. Then
\begin{align*}
 B_{ab,cd}^2-B_{ac,bd}^2&=2k(b-c)(a-b-c-d),\\
 B_{ac,bd}^2-B_{ad,bc}^2&=2k(a-b)(c-d).
\end{align*}
The dominant case, including its boundary, is strict by
Lemma~\ref{lem:dominant}. In the other case a largest norm uses the
cross pairing $(a,c),(b,d)$. The moment formula is
\[
 A=\max\{a,(a+b+c)/2,(3a+b+c+d)/4\}.
\]
Indeed, averaging the sign of $a$ yields the maximum of $a$ and the
absolute remaining sum. Only $b+c+d$ and $b+c-d$ among the four
possible remaining magnitudes can exceed $a$.

Write $(a,b,c,d)=(x+u,x-u,y+v,y-v)$ in the nondominant case.
Then $0\le u,v\le y$ and $x\ge y+u+v$. For the reference
$(x,x,y,y)$ paired as $(x,y),(x,y)$, write
$A_0=x+y/2$, $B_0=\sqrt2(x+ky)$, $Q_0=2(x^2+y^2)$.
Put $d_0=\max(u,v)$. Evaluating the second and third branches of
the moment formula gives $A-A_0\ge\max(v/2,u/2)=d_0/2$.
The squared differences are $B^2-B_0^2=2(u+kv)^2$ and
$Q-Q_0=2(u^2+v^2)\ge c_2(B^2-B_0^2)$.
For $d_0>0$ rationalization gives
\[
 B-B_0\le \frac{(u+kv)^2}{\sqrt2(x+ky)}
 \le d_0\frac{u+kv}{\sqrt2d_0+u+v}\le k d_0.
\]
The second step uses $u+kv\le\sqrt2d_0$ and
$x+ky\ge\sqrt2d_0+u+v$. The last step uses $(1-k)u\le k\sqrt2d_0$. In the deficit
$r^2+2rA+Q-\gamma(r+B)^2$, the linear term's increase is
nonnegative because $\gamma k<1/2$, and the constant term strictly
increases when $d_0>0$ because $c_2>\gamma$. When $d_0=0$ the
reference is the original sorted vector. It suffices to optimize
the reference and any equality must be symmetric.

Put $s=\sqrt2$, $z=(r,x,y)^\top$ and $\ell=(1,s,2-s)^\top$.
The reference quotient is $z^\top M z/(\ell^\top z)^2$, where
\[
 M=\begin{pmatrix}1&1&1/2\\1&2&0\\1/2&0&2\end{pmatrix}.
\]
The Schur complement of the lower block is $3/8>0$.
Cauchy--Schwarz in the $M$ inner product yields the reciprocal of
$\ell^\top M^{-1}\ell=(15-8s)/3$ as a lower bound. The vector
\[
 M^{-1}\ell=\bigl((4-2s)/3,\ (5s-4)/6,\ (2-s)/3\bigr)^\top
\]
has $r>0$ and $x>y>0$, so its positive ray attains this bound. At
$y=1$ it gives $r=2$ and $x=x_0:=(1+3s)/2$. Both maximizing
pairs must be $(x_0,1)$ because pairing the like entries strictly
reduces $B$. Positive lighter coefficients force equal leaves
on reversing the within-star replacement. The formula includes
zero padding, which cannot give equality on a missing required
leaf or star. At the witness
$\Vstar=(39+18\sqrt2)/2$ and $U=(99-14\sqrt2)/2$.
\end{proof}

Fang and Wang \cite[Theorem 1]{fangwang} prove that, for sorted
nonnegative unit vectors of dimension at least four at distance
$0\le\tau\le\sqrt{2-\sqrt2}$ from $(1,1,0,\ldots)/\sqrt2$,
a profile $(x,x,y,y,0,\ldots)$ uniquely minimizes the first absolute
Rademacher moment. The normalized equality vector in
Theorem~\ref{thm:seven} has distance about $0.36243$, within that
range. The quotient optimizes over the distance, outer weight and
matching as well. For the same coefficient profile, crossing
heavy with light gives $B^2=2(x_0+k)^2$, while pairing like entries
gives $B^2=Q$. These pairings share the first moment and distance.

\begin{theorem}[Three small child stars]\label{thm:three}
With at most three disjoint child stars and at most two active leaves
per star, the sharp factor in $\Vstar\ge\gamma U$ is the same
$\gamma$. At $U>0$, equality uses the weights of
Theorem~\ref{thm:seven} with any remaining star of weight zero at every
node. If all three stars have positive total weight, the bound is
strict. The infimum over that subclass is $\gamma$, approached on
trees with at least two child stars that each have two leaves
available.
\end{theorem}
\begin{proof}
Minimize the quadratic deficit at fixed coefficients over $r\ge0$.
Its minimizer is $r=(\gamma B-A)_+/(1-\gamma)$, so the target is
equivalent to
\begin{equation}
 (\gamma B-A)_+^2\le(1-\gamma)(Q-\gamma B^2).
 \label{eq:deficit}
\end{equation}
Here $Q\ge c_2B^2$ and $A\le\sqrt Q\le B$. The last comparison
follows pairwise from $(a+kb)^2-a^2-b^2=2kb(a-b)$. For
$H=(1-\gamma)/\gamma$ we have
\[
 (\gamma B-A)^2-(1-\gamma)(Q-\gamma B^2)
 =\gamma\bigl((B-A)^2-H(Q-A^2)\bigr).
\]
Consequently the norm inequality
\begin{equation}
 B-A\le\sqrt{H(Q-A^2)}
 \label{eq:norm-test}
\end{equation}
suffices. A strict norm inequality excludes equality.
If $\gamma B\le A$ and $B>0$, strictness follows already from
$Q>\gamma B^2$. A nonzero outer weight with $B=0$ gives ratio one.
Pad to three pairs with zero pairs where needed. Lemma~\ref{lem:dominant}
handles the dominant region. For all other coefficients,
Appendix~\ref{app:tape} proves~\eqref{eq:norm-test} using a finite
exact certificate. Its cone inequalities hold on
all nonnegative generator combinations.

The appendix also shows that equality in the nonzero case forces a
positive multiple of $(x_0,x_0,1,1,0,0)$ in sorted coefficients.
Only two heavy--unit pairs and a zero pair maximize $B$ on this
ray. The deficit strictly decreases with $B>0$ at fixed $r,A,Q$,
so all other matchings give strictness. The unique minimizing
outer weight is $2$ in units of the lighter active coefficient,
and the equality leaves follow from the within-star replacement.
The seven-weight witness attains equality. Adding positive weights
to the third star and taking them to zero gives the same infimum
by continuity of the finite formulas.
\end{proof}
